\begin{figure}[htbp]
\centering
\small

\begin{adjustbox}{max width=\linewidth}
\begin{tikzpicture}[x=1.45cm, y=1cm]

  \fill[phaseIcol]  (0,0) rectangle (3,0.65);
  \fill[phaseIIcol] (3,0) rectangle (7,0.65);
  \fill[phaseIcol]  (7,0) rectangle (11,0.65);

  \draw[gray!45, very thin] (0,0) rectangle (11,0.65);
  \foreach \x in {3,7}{
    \draw[gray!45, very thin] (\x,0) -- (\x,0.65);
  }

  \node[font=\small\bfseries, text=sigbuy]             at (1.5,0.325) {Regime~I};
  \node[font=\small\bfseries, text=SteelBlue!80!black] at (5.0,0.325) {Regime~II};
  \node[font=\small\bfseries, text=sigbuy]             at (9.0,0.325) {Regime~III};

  \foreach \m/\lbl in {
      0/Aug, 1/Sep, 2/Okt, 3/Nov, 4/Dez,
      5/Jän, 6/Feb, 7/Mär, 8/Apr, 9/Mai, 10/Jun, 11/{}}{
    \draw[gray!55, thin] (\m,-0.05) -- (\m,0);
    \node[anchor=north, font=\footnotesize, text=gray!75!black]
      at (\m,-0.08) {\lbl};
  }

  \foreach \xA/\xB/\lbl in {
      0/3/{ca. 4 Monate},
      3/7/{ca. 3 Monate},
      7/11/{ca. 4 Monate, bis Datenende}}{
    \draw[{<[scale=0.7]}-{>[scale=0.7]}, gray!60, thin]
      (\xA,-0.56)--(\xB,-0.56)
      node[midway, below, font=\footnotesize, text=gray!75!black]{\lbl};
  }

\end{tikzpicture}
\end{adjustbox}

\vspace{2.0em}

\begin{adjustbox}{max width=\linewidth}
\begin{tikzpicture}[x=1cm, y=1cm]

  \def\slotwd{0.16}
  \def\slotht{0.26}
  \def\slotstep{0.27}
  \providecommand{\tariffslot}[3]{%
    \fill[#3] (#1,#2) rectangle (#1+\slotwd,#2+\slotht);%
  }

  \def\xI{0}
  \def\xII{5.25}
  \def\xIII{10.50}

  % Regime I
  \node[font=\small\bfseries, text=sigbuy]
    at (\xI+1.75,1.72) {Regime~I};
  \node[
    font=\scriptsize,
    text=gray!70!black,
    align=center,
    text width=3.3cm
  ] at (\xI+1.75,1.22) {dünn besetzt,\\ zoneneinheitlich};

  \foreach \yrow in {0.10,0.42,0.74}{
    \foreach \s/\col in {
        0/signeutral, 1/signeutral, 2/sigbuy,
        3/signeutral, 4/signeutral, 5/signeutral,
        6/sigsell,    7/signeutral, 8/signeutral,
        9/sigbuy,    10/signeutral,11/signeutral,
       12/signeutral}{
      \tariffslot{\xI+\s*\slotstep}{\yrow}{\col}
    }
  }

  \foreach \yrow/\lbl in {0.23/{HH\,1}, 0.55/{HH\,2}, 0.87/{HH\,3}}{
    \node[anchor=east, font=\footnotesize, text=gray!70!black]
      at (\xI-0.08,\yrow) {\lbl};
  }

  \node[
    font=\tiny,
    text=gray!65!black,
    align=center,
    text width=3.3cm
  ] at (\xI+1.75,-0.20) {innerhalb der Netzzone identisch};

  % Regime II: tatsächliche stündliche Signalzustände an zwei Tagen
  \def\hourwd{0.115}
  \def\hourht{0.11}
  \def\hourstep{0.145}
  \providecommand{\hourslot}[3]{%
    \fill[#3] (#1,#2) rectangle (#1+\hourwd,#2+\hourht);%
  }
  \newcommand{\drawhourrow}[3]{%
    \foreach \state [count=\hh from 0] in #3 {
      \if\state H \def\sc{sigsell}\else \def\sc{sigbuy}\fi
      \hourslot{#1+\hh*\hourstep}{#2}{\sc}
    }%
  }

  \node[font=\small\bfseries, text=SteelBlue!80!black]
    at (\xII+1.75,1.72) {Regime~II};
  \node[
    font=\scriptsize,
    text=gray!70!black,
    align=center,
    text width=3.3cm
  ] at (\xII+1.75,1.22) {dicht,\\ Gruppenrotation};

  % 27. November 2025, oben: Rotationsgruppen 1 bis 3
  \drawhourrow{\xII}{0.775}{{L,L,L,H,H,L,L,L,L,H,H,L,L,L,L,H,H,L,L,L,L,L,H,H}}
  \drawhourrow{\xII}{0.630}{{H,L,L,L,L,H,H,L,L,L,L,H,H,L,L,L,L,H,H,L,L,L,L,L}}
  \drawhourrow{\xII}{0.485}{{L,H,H,L,L,L,L,H,H,L,L,L,L,H,H,L,L,L,L,L,H,H,L,L}}

  % 28. November 2025, unten: dieselben drei Rotationsgruppen
  \drawhourrow{\xII}{0.340}{{H,H,H,H,L,L,H,H,H,H,L,L,H,H,H,H,L,L,H,H,H,H,H,L}}
  \drawhourrow{\xII}{0.195}{{L,L,H,H,H,H,L,L,H,H,H,H,L,L,H,H,H,H,H,L,L,H,H,H}}
  \drawhourrow{\xII}{0.050}{{H,H,L,L,H,H,H,H,L,L,H,H,H,H,L,L,H,H,H,H,H,L,L,H}}

  \node[anchor=east, font=\scriptsize, text=gray!70!black]
    at (\xII-0.05,0.63) {27. Nov.};
  \node[anchor=east, font=\scriptsize, text=gray!70!black]
    at (\xII-0.05,0.195) {28. Nov.};

  \node[
    font=\tiny,
    text=gray!65!black,
    align=center,
    text width=3.3cm
  ] at (\xII+1.75,-0.20) {täglich wechselnde Blockgrenzen};

  % Regime III
  \node[font=\small\bfseries, text=sigbuy]
    at (\xIII+1.75,1.72) {Regime~III};
  \node[
    font=\scriptsize,
    text=gray!70!black,
    align=center,
    text width=3.3cm
  ] at (\xIII+1.75,1.22) {dünn besetzt,\\ zoneneinheitlich};

  \foreach \yrow in {0.10,0.42,0.74}{
    \foreach \s/\col in {
        0/signeutral,  1/sigbuy,    2/signeutral,
        3/signeutral,  4/signeutral, 5/signeutral,
        6/sigsell,     7/signeutral, 8/signeutral,
        9/signeutral, 10/sigbuy,   11/signeutral,
       12/signeutral}{
      \tariffslot{\xIII+\s*\slotstep}{\yrow}{\col}
    }
  }

  \node[
    font=\tiny,
    text=gray!65!black,
    align=center,
    text width=3.3cm
  ] at (\xIII+1.75,-0.20) {wie Regime~I aufgebaut};

  \node[font=\footnotesize, text=gray!60!black, align=center]
    at (7.05,-1.10)
    {Regime~I und III: schematische Viertelstundenintervalle,
     nicht maßstabsgetreu. \\
     Regime~II: tatsächlicher stündlicher Signalzustand,
     27.--28. November 2025.};

\end{tikzpicture}
\end{adjustbox}

\vspace{1.4em}

\begin{minipage}{0.94\linewidth}
\centering
\footnotesize
\begin{tabular}{@{}lll@{}}
\tikz{\fill[sigbuy] (0,0) rectangle (0.25,0.18);} &
Niedrigtarifsignal &
Netzbezug in dieses Intervall verschieben \\[0.35em]

\tikz{\fill[sigsell] (0,0) rectangle (0.25,0.18);} &
Hochtarifsignal &
Netzbezug in diesem Intervall reduzieren \\[0.35em]

\tikz{\fill[signeutral] (0,0) rectangle (0.25,0.18);
      \draw[gray!40] (0,0) rectangle (0.25,0.18);} &
Normaltarifintervall &
regulärer Netztarif
\end{tabular}
\end{minipage}

\vspace{1.0em}

\caption{Tarifsignalregime und Signalstruktur auf Haushaltsebene}
\caption*{\footnotesize \textit{Anmerkungen:}
  Teil~A zeigt die drei Regime entlang des Interventionszeitraums.
  Teil~B veranschaulicht die Signalstruktur für drei behandelte
  Rotationsgruppen innerhalb einer Netzzone. Die Zeilen für Regime~I
  und Regime~III sind schematische Beispiele: Beide Regime weisen eine
  dünn besetzte, prognosebasierte Struktur auf, bei der alle
  Tarifhaushalte einer Netzzone dieselbe Folge erhalten und die meisten
  Intervalle im Normaltarif liegen. Die Zeilen für Regime~II sind nicht
  schematisch, sondern zeigen die tatsächlich realisierten stündlichen
  Signalzustände der drei Rotationsgruppen an zwei aufeinanderfolgenden
  Tagen, dem 27. und 28. November 2025. In Regime~II ist jedes Intervall
  aktiv; ein Normaltarifzustand besteht nicht. Die Gruppen erhalten
  unterschiedliche Signalfolgen. Anders als in Regime~I und III
  verlaufen die Signale in Regime~II in mehrstündigen Blöcken, deren
  Grenzen sich von einem Kalendertag zum nächsten verschieben. Die
  Zuordnung zu den Rotationsgruppen ist unabhängig vom Haushaltsverbrauch.
  Die Tarifstufen sind in allen drei Regimen identisch; nur Signaldichte,
  Vorhersagbarkeit und Rotationsstruktur ändern sich.}
\label{fig:nooe_tariff_regimes}
\end{figure}
